%! Carrying Out a Test for the Difference of Two Population Proportions — Unit 6, Lesson 6.11 — Homework
\documentclass[10pt]{article}
\usepackage{apstats-article}
\usepackage{apstats-boxes}
\newcommand{\TallMath}[1]{$\displaystyle #1\rule[-1.4em]{0pt}{3.2em}$}

\begin{document}

\pageheader{Unit 6, Lesson 6.11}{Homework}
\namedateperiod

\vspace{0.15in}

Complete the full four-step procedure for each problem. Use $\alpha = 0.05$ unless stated otherwise.

\begin{enumerate}

\item \textbf{Social Media and Academic Performance.} A researcher hypothesizes that
students who use social media more than 2 hours per day have a higher proportion of
low-GPA outcomes than those who use it 2 hours or less per day.
High use: 54 of 180 had a low GPA.
Low use: 36 of 200 had a low GPA.

\begin{enumerate}[label=\alph*.]
  \item State $H_0$ and $H_a$ (one-sided). \writelines{2}
  \item Compute $\hat p_c$ and verify conditions. \writelines{3}
  \item Compute $z$ and the $p$-value. \writelines{3}
  \item Conclude in context. \writelines{2}
\end{enumerate}

\item \textbf{Employee Turnover.} A company compares turnover rates for employees who
went through a new onboarding program versus the old one.
New program: 18 of 240 left within the first year.
Old program: 30 of 240 left within the first year.
Test whether the new program reduced turnover (one-sided, $\alpha=0.05$).

\begin{enumerate}[label=\alph*.]
  \item State $H_0$ and $H_a$. \writelines{2}
  \item Compute $\hat p_c$ and verify conditions. \writelines{3}
  \item Compute $z$ and the $p$-value. \writelines{3}
  \item Conclude in context. \writelines{2}
\end{enumerate}

\end{enumerate}

\end{document}
